\chapter*{Resumen}
\addcontentsline{toc}{chapter}{Resumen}
\markboth{\MakeUppercase{Resumen}}{\MakeUppercase{Resumen}}

El presente trabajo tuvo por objetivo desarrollar una biblioteca de Python que permitiera ejecutar operaciones tensoriales e inferencia de redes neuronales sobre datos cifrados con el esquema de cifrado homomórfico CKKS, aprovechando la capacidad de procesamiento paralelo del GPU. Para resolverlo se diseñó una interfaz declarativa de alto nivel, análoga a la de las bibliotecas de aprendizaje profundo de uso común. Se buscó ocultar del usuario lo más posible la complejidad del esquema criptográfico y la gestión de la memoria de video. Se implementaron las capas lineal, convolucional y de activación traduciendo cada operación de la red a las primitivas homomórficas del esquema, junto con una calibración automática del modelo y una gestión automática del presupuesto de ruido que habilitó el \textit{bootstrapping} cuando la profundidad del modelo lo exigió. El cómputo criptográfico se delegó a HEonGPU, una biblioteca especializada que ejecuta las primitivas de CKKS sobre la GPU, y el desarrollo se guió por un conjunto de pruebas automatizadas. La biblioteca se evaluó sobre dos casos de uso representativos, una red totalmente conectada para regresión y una red convolucional para clasificación, comparándola con una ejecución de referencia en texto plano y con dos bibliotecas del estado del arte en términos de tiempo, consumo de memoria y fidelidad de los resultados; la inferencia cifrada reprodujo el cómputo en texto plano dentro del margen de error propio de la aritmética aproximada de CKKS. La biblioteca resultante fue publicada como paquete instalable de código abierto y disponible para la comunidad, que pone la inferencia privada acelerada por GPU al alcance de usuarios sin conocimientos especializados en criptografía.

\vspace{2em}

\noindent\textbf{Palabras clave:} cifrado homomórfico, CKKS, inferencia privada, redes neuronales, aprendizaje automático, GPU, CUDA, privacidad, Python, HEonGPU.
